\chapter{渎职罪}

% Source PDF physical page 427; see audit/source_coverage.json.
\begin{EditionOpeningPanel}
重点罪名是玩忽职守罪、徇私枉法罪。
\end{EditionOpeningPanel}

\section{重点罪名}

[问题]巡警甲在大街上看到乙在杀丙，袖手旁观。乙顺利杀死了丙。甲构成滥用职权罪还是玩忽职守罪？（袖手旁观案）

\subsection*{一、滥用职权罪和玩忽职守罪（第397条）}

\subsubsection*{（一）滥用职权罪}

滥用职权罪，是指国家机关工作人员滥用职权，致使公共财产、国家和人民利益遭受重大损失的行为。

1.行为方式。

（1）擅权。这是指故意不正确履行职责。

（2）弃权。这是指故意不履行应当履行的职责。这表明，滥用职权罪可以由不作为构成。（2023年试题）

（3）越权。这是指超越职权处理事项。

2.实害结果。成立本罪要求致使公共财产、国家和人民利益遭受重大损失。注意：该实害结果不要求与滥用职权行为具有直接的因果关系。

3.主观必须是故意。注意：对于不要求有直接因果关系的实害结果，不要求行为人认识到，只要有认识可能性即可。

4.罪数。司法解释规定：国家工作人员不以个人占有为目的，利用职务上的便利，用公共财物向他人行贿，同时构成行贿罪和渎职犯罪的，数罪并罚。\footnote{2026年5月1日最高人民法院、最高人民检察院《关于办理贪污贿赂刑事案件适用法律若干问题的解释

（二）》。}该数罪并罚存在特殊规定的情形。

% Source PDF physical page 428; see audit/source_coverage.json.
例1，官员甲为了给本单位谋取不正当利益，未经集体决定，私下将公款从公账直接打到官员乙账户。甲不构成贪污罪，因为是为单位办事，没有个人占有公款的目的。甲不构成单位行贿罪，而构成行贿罪。甲也构成滥用职权罪。甲的一个行为触犯两罪，本应按照想象竞合，择一重罪论处，但是，该司法解释要求数罪并罚。这属于罪数上的特殊规定，不能推广适用。

例2，官员甲为了给自已谋取不正当利益，带着不归还的意图，私下将公款从公账直接打到官员乙账户。甲构成贪污罪，因为为了给自己办事，将公款白白送人，对公款具有非法占有目的。甲还构成行贿罪、滥用职权罪。贪污罪、行贿罪、滥用职权罪，按照想象竞合，择一重罪论处。由于本案前提事实不符合上述司法解释的前提事实，所以不需要遵守该司法解释的特殊规定（数罪并罚）。

\subsubsection*{（二）玩忽职守罪}

玩忽职守罪，是指国家机关工作人员玩忽职守，致使公共财产、国家和人民利益遭受重大损失的行为。

1.行为方式：一是不履行职责，即不作为的方式；二是不正确履行职责。

2.主观必须是过失。

[区分］滥用职权罪和玩忽职守罪的区分：滥用职权罪是故意犯罪，故意实施违背职责的行为；玩忽职守罪是过失犯罪，过失实施违背职责的行为，也即不认真履行职责。

例如（2019年试题），夜里，甲在大街上欲杀害乙。乙打了几次报警电话，说有人杀自己。由于乙当时有点醉酒，口齿不是很清楚，警察以为乙是恶作剧，没有出警。乙被甲杀死。如果警察及时出警，乙不会被杀死。警察构成玩忽职守罪。

[提示]认识上容易犯的错误是：将不作为的滥用职权罪当作玩忽职守罪。例如，本节开头“袖手旁观案”中，甲的这种行为不是玩忽职守，而是滥用职权，构成滥用职权罪，也构成故意杀人罪的不作为的帮助犯，想象竞合，择一重罪论处。

[对比练习]例1，警察丁利用职权，收受贿赂，明知田某等人没有驾驶技术，不能获取驾驶证，仍发放驾驶证。丁多次实施这种渎职行为。某日，田某由于没有驾驶技术，驾车时酿成车祸，导致多人死亡。丁构成滥用职权罪。第一，成立滥用职权罪，在客观上不要求该车祸与丁的滥用职权行为有直接的因果关系。第二，主观上不要求丁认识到该车祸发生，只要有认识到的可能性即可。（2021年试题）

例2（2025年试题），甲是车管所的体检工作人员，为了图省事，未对申请人乙做体检，便出具视力合格的报告，致使存在视力缺陷的乙领取大型货车驾驶证。此后半年，乙因视力缺陷发生车祸，造成多人死亡。甲的行为和事故之间有因果关系。

[对比]主体身份不同，罪名不同。刑法第168条规定：国有公司、企业、事业单位人员由于严重不负责任或者滥用职权，致使国家利益遭受重大损失，构成国有公司、企业、事业单位人员失职罪或者国有公司、企业、事业单位人员滥用职权罪。

[注意]单位与个人的问题。以“集体研究”形式实施的渎职犯罪，应当\CJKunderwave{追究国家机}\CJKunderwave{关负有责任的人员的刑事责任}。

% Source PDF physical page 429; see audit/source_coverage.json.
\subsection*{二、徇私枉法罪和民事、行政枉法裁判罪（第399条）}

\subsubsection*{（一）徇私枉法罪}

徇私枉法罪，是指司法工作人员私枉法、徇情枉法，对明知是无罪的人而使他受追诉，对明知是有罪的人而故意包庇不使他受追诉，或者在刑事审判活动中故意违背事实和法律作枉法裁判的行为。

1.行为主体：司法工作人员，具体而言是，刑事案件的侦查人员（包括勘验、检查人员）、检察人员、审判人员。

（1）监察机关中负责调查职务犯罪的人员，属于本罪的主体。

（2）海关中负责侦查走私犯罪的人员，属于本罪的主体。

[注意］本罪行为主体不是一切司法工作人员，仅限于参与案件追诉、审理的司法工作人员。例如，王法官承办甲的抢劫案件。与本案无关的某派出所所长李某利用职务便利，制造虚假贩毒案件，让甲揭发，导致甲构成立功。李某不构成徇私枉法罪，而构成滥用职权罪和帮助犯罪分子逃避处罚罪，想象竞合。

2.法官徇私枉法的表现：在定罪、量刑上滥用职权。

3.检察官询私枉法的表现：在批捕、审查起诉上滥用职权。

4.侦查人员徇私枉法的表现：在立案、采取强制措施、调查证据上滥用职权。（2009年第65题）

5.主观要件是故意，且要求有徇私动机。

（1）规定徇私动机，旨在将由于法律水平不高而造成的错误排除在外。例如，法官甲法律水平很高，知道亲戚乙无罪，但是为了彰显自已大公无私、公正廉洁，故意判乙有罪。甲实际上仍有徇私动机，这里“私”就是博取名声。甲构成徇私枉法罪。（2023年试题）

（2）行为人法律水平是合格的，但是由于过失，导致无罪的人受追诉，或者有罪的人未被追诉的，不成立徇私枉法罪，但构成玩忽职守罪。例如，法官甲法律水平挺高，但是在审理某起案件时，未认真阅卷，疏忽大意，未注意到某项证据的存在，由此认为被告人李某有罪，判处六个月有期徒刑，并予以执行。后查证李某的行为不构成犯罪。甲构成玩忽职守罪。（2025年试题）

（3）行为人由于法律水平低下，导致无罪的人受追诉，或者有罪的人未被追诉的，不构成玩忽职守罪，因为行为人不存在犯罪过失。

\subsubsection*{（二）民事、行政枉法裁判罪}

这是指司法工作人员在民事、行政审判活动中故意违背事实和法律作枉法裁判，情节严重的行为。

[注意］与受贿罪的罪数关系。

第399条第4款：“司法工作人员收受贿赂，有前三款行为的（徇私枉法罪，民事、行政枉法裁判罪，执行判决、裁定失职罪，执行判决、裁定滥用职权罪），同时又构成受贿罪的，依照处罚较重的规定定罪处罚。”这表明，除此款规定之外，国家工作人员犯其他罪同时受贿的，一律数罪并罚。（2013年第21题，2017年第62题）

% Source PDF physical page 430; see audit/source_coverage.json.
\section{普通罪名}

\subsection*{一、私放在押人员罪（第400条）}

这是指，\CJKunderwave{司法工作人员}私放\CJKunderwave{在押}的\CJKunderwave{犯罪嫌疑人、被告人或者罪犯}的行为。

1.行为对象：依法被关押的犯罪嫌疑人、被告人、罪犯，不包括被行政拘留、司法拘留的人。

2.主观是故意。对应的过失犯罪叫“失职致使在押人员脱逃罪”。（2010年第2题）

3.罪数。例如（2011年第56题），狱警甲明知囚犯乙脱逃后会杀害证人，仍然私放。乙果然杀了证人。甲同时触犯私放在押人员罪和故意杀人罪（帮助犯），想象竞合，择一重罪论处。

\subsection*{二、徇私舞弊不移交刑事案件罪（第402条）}

这是指，行政执法人员徇私舞弊，对应当移交司法机关追究刑事责任的不移交，情节严重的行为。

1.行为主体

行为主体是行政执法人员，主要是市场监督管理、税务、监察等行政执法人员。例如，环保部门的执法人员在执法检查时发现甲实施污染环境罪，故意不移交刑事案件侦查机关，构成本罪。（2023年试题）

[注意]警察的不同身份。

（1）警察作为行政执法人员。

例1，警察在处理甲的治安违法行为（小偷小摸）时，发现甲实际上偷了一万元，构成盗窃罪，但是故意不移交给刑事案件侦查部门，构成本罪。

例2，警察在治安执法时，发现甲在拐卖妇女、儿童，故意不移交给刑事案件侦查部门，构成本罪。

（2）警察作为刑事侦查人员。例如，警察在侦查甲的盗窃罪时，应当采取刑事强制措施，故意不采取，构成徇私枉法罪。

2.行为方式

本罪是不作为犯，也即不履行移交义务。不履行的方式既可以是积极举动，也可以是消极静止。积极举动有：第一，以罚代刑；第二，毁灭、隐藏证据，伪造证据，改变刑事案件性质（同时构成帮助毁灭、伪造证据罪，想象竞合）。

3.主观要件是故意，且须有徇私动机。（2023年试题）

4.想象竞合

本罪与放纵走私罪（第411条）、放纵制售伪劣商品犯罪行为罪（第414条）、帮助犯罪分子逃避处罚罪（第417条），可以想象竞合。

\subsection*{三、帮助犯罪分子逃避处罚罪（第417条）}

这是指，有查禁犯罪活动职责的国家机关工作人员，向犯罪分子通风报信、提供便利，帮助犯罪分子逃避处罚的行为。

1.帮助逃避处罚的情形：帮助逃避侦查、追诉、定罪、重罪重刑处罚、刑罚执行。

% Source PDF physical page 431; see audit/source_coverage.json.
也即，帮助逃避受到刑事处罚。

2.行为主体与行为方式

（1）司法机关中，承办刑事案件的司法工作人员，利用职权获得办案消息，向犯罪分子通风报信，帮助其逃避处罚。该行为构成本罪，也构成徇私枉法罪，想象竞合，择一重罪论处。有些行为构成徇私枉法罪，但不构成本罪，例如，法官将有罪判为无罪。所以，两罪是交叉重合关系。

（2）司法机关中，不承办该刑事案件，但有查禁犯罪活动职责的司法工作人员，利用职权获得办案消息，向犯罪分子通风报信，帮助其逃避处罚。例如，甲犯抢劫罪，由警官乙负责侦办。第一，丙是该公安局管理网络的警察，利用职权，获知乙的侦查动向，向甲通风报信。丙构成本罪，但不构成徇私枉法罪，因为不是该案的办案人员。第二，丁是该公安局的档案室管理员，偶尔听到乙的侦查动向，向甲通风报信。丁不负有查禁犯罪活动的职责，不构成帮助犯罪分子逃避处罚罪，构成窝藏罪。

（3）其他国家机关中，有查禁犯罪活动职责的国家机关工作人员，利用职权，帮助犯罪分子逃避处罚。例如，某地农业农村局的执法人员甲负有查禁破坏鱼类资源犯罪活动的职责，发现乙实施非法捕捞犯罪。第一，甲不予处理，让乙赶紧逃。甲构成帮助犯罪分子逃避处罚罪。第二，甲处理后，故意不向负责刑事案件的司法机关移交。甲构成徇私舞弊不移交刑事案件罪。

3.罪数。本罪可以和其他罪想象竞合。

（1）利用职权，为犯罪分子通风报信、提供便利条件，帮助逃匿，构成本罪，同时构成窝藏罪，想象竞合。

（2）帮助、示意犯罪分子隐匿、毁灭、伪造证据，逃避处罚，构成本罪，同时构成帮助毁灭、伪造证据罪，想象竞合；帮助、示意犯罪分子串供、翻供，构成本罪，同时构成妨害作证罪，想象竞合。

（3）公安人员对盗窃、抢劫的机动车辆，非法提供机动车辆牌证，或者为其取得机动车辆牌证提供便利，构成本罪，同时构成掩饰、隐瞒犯罪所得罪，想象竞合。（2017年第63题）

\subsection*{四、不解救被拐卖、绑架妇女、儿童罪（第416条）}

这是指，负有解救职责的国家机关工作人员接到被拐卖、绑架的妇女、儿童及其亲属的解救要求，或者接到其他人的举报，不进行解救，造成严重后果的行为。

1.行为对象是被拐卖、绑架的妇女、儿童。（1）其中的“绑架”是指绑架型拐卖罪，不是指绑架罪。（2）被拐卖的妇女、儿童，既包括尚未卖掉的妇女、儿童，也包括已经卖掉的妇女、儿童。

2.对非法拘禁罪、绑架罪、拐骗儿童罪的被害人，有解救职责，故意不解救，定滥用职权罪；过失不解救，定玩忽职守罪。（2023年试题）

\subsection*{五、食品、药品监管渎职罪（第408条之一）}

该罪是指，负有食品药品安全监督管理职责的国家机关工作人员，滥用职权或者玩忽职守，造成严重后果或者有其他严重情节的行为。\footnote{该罪名实际上应当分解为两个罪名：食品、药品监管滥用职权罪和食品、药品监管失职罪。例如，刑法第397条这一个条文中包括滥用职权和玩忽职守，但被分解为两个罪名：滥用职权罪和玩忽职守罪。}

